\documentclass[11pt]{article}

\usepackage[margin=1in]{geometry}
\usepackage{amsmath,amssymb,amsthm}
\usepackage{graphicx}
\usepackage[hidelinks]{hyperref}
\emergencystretch=3em

\newtheorem{theorem}{Theorem}
\newtheorem{lemma}{Lemma}

\newcommand{\N}{\mathbb N}
\newcommand{\R}{\mathbb R}
\newcommand{\M}{\mathcal M}
\newcommand{\X}{\mathcal X}
\newcommand{\Lip}{\operatorname{Lip}}
\newcommand{\linspan}{\operatorname{span}}

\title{A Compact \(c_0\) Subset With No Lipschitz \(p\)-Approximate Schauder Frame}
\author{Run fa\_banach\_001, lane 7}
\date{June 16, 2026}

\begin{document}
\maketitle

\section*{Status}

This packet gives a counterexample to the universal part of Problem~3.1(1) in
Krishna--Johnson, arXiv:2211.10652 \cite{KJ2022}. Not every subset of a
Banach space has a Lipschitz \(p\)-approximate Schauder frame for some
\(1\leq p<\infty\). In fact, there is a compact subset of a Banach space
isometric to \(c_0\) which has no Lipschitz \(p\)-ASF for any finite \(p\).

\section*{Source Problem}

The source paper formulates the following open problem on PDF page~10.

\begin{figure}[h]
\centering
\includegraphics[width=\textwidth]{figures/open_problem_crop.png}
\caption{Problem~3.1 from arXiv:2211.10652, PDF page~10.}
\end{figure}

The construction below answers the ``in particular'' question in item~(1)
negatively. It does not claim a classification of all subsets and does not
settle item~(2)'s weaker representation problem.

\section*{Definitions Used}

For a subset \(\M\subseteq \X\), a Lipschitz \(p\)-ASF in the sense of
\cite{KJ2022} consists of scalar coefficient maps \(f_n\), points
\(\tau_n\in\M\), an analysis map
\[
  \theta_f:\M\to \ell^p,\qquad \theta_f(x)=(f_n(x))_{n=1}^\infty,
\]
a bounded synthesis operator
\[
  \theta_\tau:\ell^p\to \X,\qquad
  \theta_\tau((a_n))=\sum_{n=1}^\infty a_n\tau_n,
\]
and a Lipschitz frame map
\[
  S_{f,\tau}=\theta_\tau\theta_f:\M\to\M
\]
which is invertible and bi-Lipschitz.

\section*{Proof Intuition}

The analysis map of a Lipschitz \(p\)-ASF is not merely a coefficient map: it
must bi-Lipschitzly encode the metric of \(\M\) into \(\ell^p\). Indeed, the
frame map is the synthesis operator applied to the analysis map, and the frame
map has a Lipschitz inverse. Thus a subset with no bi-Lipschitz embedding into
any \(\ell^p\), \(p<\infty\), cannot have a Lipschitz \(p\)-ASF.

The counterexample is a compact bouquet in \(c_0\) containing smaller and
smaller copies of the balls of \(\ell_\infty^n\) for all \(n\). Compactness
comes from shrinking the \(n\)-dimensional pieces to zero. The obstruction is
local: a bi-Lipschitz embedding into a fixed \(\ell^p\) would, on every
finite-dimensional piece, produce uniformly distorted embeddings of
\(\ell_\infty^n\) into \(\ell^p\). Rademacher differentiation turns these into
linear embeddings, while cotype of \(\ell^p\) forces the distortion to grow
with \(n\).

\begin{lemma}\label{lem:asf-embeds}
If \(\M\subseteq\X\) has a Lipschitz \(p\)-ASF for some \(1\leq p<\infty\),
then \(\M\) admits a bi-Lipschitz embedding into \(\ell^p\).
\end{lemma}

\begin{proof}
Let \(A=\theta_f:\M\to\ell^p\), \(V=\theta_\tau:\ell^p\to\X\), and
\(S=VA:\M\to\M\). By definition, \(A\) is Lipschitz, \(V\) is bounded, and
\(S\) is bi-Lipschitz. Thus there are constants \(L,m>0\) such that
\[
  \|A x-A y\|_p\leq L\|x-y\|,\qquad
  \|Sx-Sy\|\geq m\|x-y\|
\]
for all \(x,y\in\M\). Since \(S=VA\),
\[
  \|A x-A y\|_p
  \geq \frac{1}{\|V\|}\|V(Ax-Ay)\|
  = \frac{1}{\|V\|}\|Sx-Sy\|
  \geq \frac{m}{\|V\|}\|x-y\|.
\]
Hence \(A\) is a bi-Lipschitz embedding of \(\M\) into \(\ell^p\).
\end{proof}

\begin{lemma}\label{lem:no-lp}
There is a compact set \(\M\subset c_0\) which does not admit a bi-Lipschitz
embedding into \(\ell^p\) for any \(1\leq p<\infty\).
\end{lemma}

\begin{proof}
Let \(E_n=\ell_\infty^n\) and let
\[
  X_0=\left(\bigoplus_{n=1}^{\infty}E_n\right)_{c_0}.
\]
This Banach space is isometric to \(c_0\), after enumerating the finitely many
coordinates in each block. Put \(r_n=2^{-n}\). We regard \(E_n\) as the
\(n\)-th coordinate block of \(X_0\), and define
\[
  \M=\{0\}\cup \bigcup_{n=1}^{\infty} r_n B_{E_n},
\]
where \(r_n B_{E_n}\) is placed in the \(n\)-th block and all other blocks are
zero.

The set \(\M\) is compact. Indeed, for every \(\varepsilon>0\) all but
finitely many blocks lie in the \(\varepsilon\)-ball around \(0\), and the
remaining union is finite-dimensional compact.

Suppose, towards a contradiction, that \(F:\M\to\ell^p\) is bi-Lipschitz for
some \(1\leq p<\infty\). Then there are constants \(0<a\leq b<\infty\) such
that
\[
  a\|x-y\|_{X_0}\leq \|F(x)-F(y)\|_p\leq b\|x-y\|_{X_0}
  \qquad (x,y\in\M).
\]
Fix \(n\). On the relative interior of the block \(r_n B_{E_n}\), define
\[
  \Phi_n(u)=r_n^{-1}F(r_n u),\qquad u\in (-1,1)^n\subset \ell_\infty^n.
\]
Then
\[
  a\|u-v\|_\infty\leq \|\Phi_n(u)-\Phi_n(v)\|_p
  \leq b\|u-v\|_\infty
  \qquad (u,v\in (-1,1)^n).
\]

Since \(\ell^p\) has the Radon--Nikodym property for \(1\leq p<\infty\),
the Banach-valued Rademacher theorem gives a point of Fréchet
differentiability of \(\Phi_n\). Let
\[
  T_n:\ell_\infty^n\to\ell^p
\]
be the derivative at such a point. Passing to difference quotients in the
upper and lower Lipschitz inequalities yields
\[
  a\|h\|_\infty\leq \|T_n h\|_p\leq b\|h\|_\infty
  \qquad (h\in\ell_\infty^n).
\]

Now let \(q=\max\{p,2\}\). The space \(\ell^p\) has finite cotype \(q\);
write \(C_p\) for a cotype \(q\) constant. Applying the cotype inequality in
\(\ell^p\) to the vectors \(T_n e_1,\dots,T_n e_n\), where \(e_i\) is the
canonical basis of \(\ell_\infty^n\), gives
\[
  \left(\sum_{i=1}^n \|T_n e_i\|_p^q\right)^{1/q}
  \leq
  C_p\,\mathbb E_\varepsilon
  \left\|\sum_{i=1}^n \varepsilon_i T_n e_i\right\|_p .
\]
The lower estimate on \(T_n\) gives the left side at least \(a n^{1/q}\).
The upper estimate gives the right side at most
\[
  C_p b\,\mathbb E_\varepsilon
  \left\|\sum_{i=1}^n \varepsilon_i e_i\right\|_\infty
  = C_p b.
\]
Thus
\[
  \frac{b}{a}\geq \frac{n^{1/q}}{C_p}
\]
for every \(n\), an impossibility. Therefore \(\M\) cannot bi-Lipschitz embed
into any \(\ell^p\), \(1\leq p<\infty\).
\end{proof}

\begin{theorem}
There is a compact subset \(\M\) of a Banach space isometric to \(c_0\) such
that \(\M\) has no Lipschitz \(p\)-approximate Schauder frame for any
\(1\leq p<\infty\).
\end{theorem}

\begin{proof}
Take the compact set \(\M\subset c_0\) from Lemma~\ref{lem:no-lp}. If \(\M\)
had a Lipschitz \(p\)-ASF for some \(1\leq p<\infty\), then
Lemma~\ref{lem:asf-embeds} would give a bi-Lipschitz embedding of \(\M\) into
\(\ell^p\), contradicting Lemma~\ref{lem:no-lp}.
\end{proof}

\section*{Verification Notes}

The only external geometry used is standard: \(\ell^p\) has cotype
\(\max\{p,2\}\), and Lipschitz maps from finite-dimensional domains into
Banach spaces with the Radon--Nikodym property are differentiable almost
everywhere. The proof does not require finite-dimensional pieces to be
linearly complemented in the target; differentiating a hypothetical
bi-Lipschitz embedding supplies the linear maps \(T_n\).

The compactness check is also important. The pieces \(r_nB_{\ell_\infty^n}\)
are infinite compact sets, not finite nets, but their radii tend to zero.
Thus every sequence either has infinitely many terms in a finite union of
compact blocks or eventually lies arbitrarily close to \(0\).

\section*{Novelty and Search Bounds}

On June~16, 2026, the run indexes were searched for arXiv:2211.10652,
``Lipschitz p-Approximate Schauder Frames'', ``Lipschitz p-ASF'',
``approximate Schauder frame'', and variants. No prior packet or attempt for
this source was found.

External web/arXiv searches for the exact title, ``Lipschitz p-ASF every
subset'', ``Classify which subsets Lipschitz p-ASF'', and ``Lipschitz
approximate Schauder frame subset Banach'' found the source paper and adjacent
frame papers by the same authors, but no later solution to Problem~3.1 and no
counterexample of the compact \(c_0\) form above. The packet should therefore
be reviewed as a new counterexample to the universal part of Problem~3.1(1),
subject to human verification.

\section*{Limitations}

This does not classify all subsets admitting Lipschitz \(p\)-ASFs. It also
does not settle Problem~3.1(2), whose stated representation property is weaker
than the Lipschitz \(p\)-ASF structure used in Lemma~\ref{lem:asf-embeds}.

\section*{Human Review Recommendation}

Review as a counterexample packet. The main points to check are:
\begin{enumerate}
\item the implication from Lipschitz \(p\)-ASF to bi-Lipschitz embeddability
into \(\ell^p\);
\item the compactness of the block bouquet \(\M\subset c_0\);
\item the differentiation/cotype obstruction for uniformly embedding all
\(\ell_\infty^n\) balls into one \(\ell^p\).
\end{enumerate}

\begin{thebibliography}{9}

\bibitem{KJ2022}
K. Mahesh Krishna and P. Sam Johnson,
\emph{Lipschitz \(p\)-Approximate Schauder Frames},
arXiv:2211.10652.

\bibitem{BL2000}
Y. Benyamini and J. Lindenstrauss,
\emph{Geometric Nonlinear Functional Analysis, Volume 1},
American Mathematical Society Colloquium Publications, vol.~48, 2000.

\bibitem{LT1991}
M. Ledoux and M. Talagrand,
\emph{Probability in Banach Spaces},
Springer, 1991.

\end{thebibliography}

\end{document}
